% REVIEW: Write Section
\section{Calls}\label{core:sec:calls}

\subsection{From the Discards}\label{core:ssec:discards}

\subsubsection{Kuikae}\label{core:sssec:Kuikae}
% REVIEW: Write section

Kuikae, or swap calling, is a rule which restricts what you may discard after making a call from the discards. It is based off of two fundamental ideals: You can't swap a tile in your hand for a copy of that tile in the discards, and you can't shift a run by calling to shift it down. These in turn lead to the following rules:
\begin{enumerate}
	\item When you make a call from the discards, the tile you choose from your hand to discard may \textit{not} be the same tile you just called.
	\item When you call Chii~\pr{call}{Chii}, you may not discard a tile on the other end of the run you Chii-ed. For example, If you Chii {\mahjong{3'45s}} you can not discard {\mahjong{6s}}.
\end{enumerate}

In the rare case where after making a call, all of your discards would be illegal by kuikae, the call made this way is considered to be illegal. Take this example hand: {\mahjong{1123m-555'z-66'6z-7'77z}}. If you call Chii on either {\mahjong{1m}} or {\mahjong{4m}} you would have no legal discards, and the call is illegal.

\call{Chii}{Chii}{チー}
	{\symbalert May only be done from the player to your left\\
	\upgradesto Must obey Kuikae~\pr{sssec}{Kuikae}\\
	\symbnegate Superceded by Pon \pr{call}{Pon}, Kan \pr{call}{Kan}, and Ron \pr{decl}{Ron}\\
	\symbnegate May not be done while there are no tiles left in the live wall}
	{As the player to your left discards a tile, you may say ``Chii,'' reveal two tiles in your hand that form a run with the discarded tile, meld them, then discard a tile. See below for examples of Chiis.
	\begin{center}These look weird, but are right: {\mahjong{3'24m}\hspace{0.2cm}\mahjong{4'23p}\hspace{0.2cm}\mahjong{2'34s}}\end{center}}

\call{Pon}{Pon}{ポン}
	{\symbnegate Superceded by Ron \pr{decl}{Ron}\\
	\upgradesto Must obey Kuikae~\pr{sssec}{Kuikae}\\
	\symbnegate May not be done while there are no tiles left in the live wall\\
	\upgradesto Upgrades into Kakan \pr{call}{Kakan}}
	{As any player discards a tile, you may say ``Pon,'' reveal two tiles in your hand that form a set with the discarded tile, meld them, then discard a tile. Please ensure the tile you grab and rotated sideways is placed correctly for where you got it from. See examples of Pons below.
	\begin{center}{From the player on your left: \mahjong{5'55m}\\ From the player across from you: \mahjong{05'5p}\\ From the player to your right: \mahjong{550's}}\end{center}}

\call{Kan}{Kan}{\ruby{明}{みん}\ruby{槓}{カン}}
	{\symbnegate Superceded by Ron \pr{decl}{Ron}\\
	\symbnegate May not be done while there are no tiles left in the live wall\\
	\upgradesto Must obey Kuikae~\pr{sssec}{Kuikae}\\
	\upgradestoother Enables Rinshan Kaihou \pr{yaku}{Rinshan Kaihou}}
	{As any player discards a tile, you may say ``Kan,'' reveal three tiles in your hand that form a set with the discarded tile, meld them, draw a tile from the dead wall, reveal a new Dora indicator, then discard a tile. Please ensure the tile you grab and rotated sideways is placed correctly for where you got it from. See examples of Kans below.
	\begin{center}{From the player on your left: \mahjong{5'505m}\\ From the player across from you: \mahjong{05'55p}/\mahjong{055'5p}\\ From the player to your right: \mahjong{5550's}}\end{center}}

\subsection{After Drawing}\label{core:ssec:drawn}

However, there are three types of Kans. The other two are able to be performed only on your turn after drawing a tile.

\call{Kakan}{Upgraded Kan}{\ruby{加}{か}\ruby{槓}{カン}}
	{\symbnegate May not be done while there are no tiles left in the live wall\\
	\upgradesfrom Upgrades from Pon \pr{call}{Pon}\\
	\upgradestoother Enables Rinshan Kaihou \pr{yaku}{Rinshan Kaihou}\\
	\upgradestoother Enables Chankan \pr{yaku}{Chankan}}
	{After you have drawn a tile, if you have a Pon of some tile, and have the fourth tile of that type in your hand, you may say ``Kan.'' Reveal that tile in your hand, place it sideways above the other sideways tile melding it, draw a tile from the dead wall, reveal a new Dora indicator, then discard a tile.
	\begin{center}{From the player on your left: \mahjong{5'05m}$\;\rightarrow\;$\mahjong{5"05m}\\ From the player across from you: \mahjong{05'5p}$\;\rightarrow\;$\mahjong{05"5p}\\ From the player to your right: \mahjong{505's}$\;\rightarrow\;$\mahjong{505"s}}\end{center}}

\call{Ankan}{Closed Kan}{\ruby{暗}{あん}\ruby{槓}{カン}}
	{\symbalert Does not open your hand\\
	\symbnegate May not be done while there are no tiles left in the live wall\\
	\upgradestoother Enables Rinshan Kaihou \pr{yaku}{Rinshan Kaihou}\\
	\upgradestoother Enables Chankan for Thirteen Orphans only \pr{yaku}{Chankan}}
	{After you have drawn a tile, if you have four copies of a tile in your hand, you may say ``Kan.'' Reveal all of those tiles, flip the outer two face down, melding it, draw a tile from the dead wall, reveal a new Dora indicator, then discard a tile.
	\begin{center}{\mahjong{6666z}$\;\rightarrow\;$\mahjong{X66Xz}}\end{center}}